\documentclass[11pt]{article}
\usepackage[margin=1.05in]{geometry}
\usepackage{fontspec}
\setmainfont{Latin Modern Roman}
\usepackage[hidelinks]{hyperref}
\usepackage{fancyhdr}
\setlength{\parskip}{.48em}
\setlength{\parindent}{0pt}
\pagestyle{fancy}
\fancyhf{}
\fancyhead[L]{\small WORKING PAPER}
\fancyhead[R]{\small Worldtext \textbullet{} 3 November 2026}
\fancyfoot[L]{\small worldtext-2026-11-03 \textbullet{} 000\_\_RETURN\_PATH.txt}
\fancyfoot[R]{\thepage}
\setlength{\headheight}{15pt}
\title{\Huge What the Shield Never Had to Decide\\[.3em]\large From Ekphrasis to an Executable Hearing}
\author{Watson Hartsoe \and Codex (drafting and compilation)}
\date{Working Paper \textbullet{} AI-augmented research process}
\begin{document}
\maketitle
\begin{abstract}
A walkable adaptation of Achilles' shield must decide how the reader can approach its hearing, who can speak, and what follows from what is said. Homer describes none of these as controls. I argue that the distinctive addition of a world model is a set of procedural commitments, and that fidelity depends in part on whether a reader can inspect and challenge them. The hearing's rival English translations, its gold, its staves, and the shield's rim make four different commitments visible. The resulting proposal treats Worldtext as a medium in which a reader can trace how an interpretation became a world that acts back.
\end{abstract}
\textbf{Keywords} ekphrasis, imagetext, world models, procedural authority, adaptation, provenance

\section{The dispute before the verdict}
An elder in a simulation of Achilles' shield might ask to see a receipt. That is a good question if the hearing concerns whether compensation was paid. It is a poor way to settle whether the claimant had the right to settle by payment in the first place. Murray's translation has the man assert that he already paid, while Butler's rendering gives him a claim about his right to pay. The difference arises before a verdict is possible. It determines what a witness or document could settle \cite{homerIliadMurray,homerIliadButler}.

A production bible that summarizes the passage as “the men dispute whether payment occurred” thereby decides what kind of proceeding a writer will stage. An interactive adaptation extends the decision further. If the reader can hand over the receipt and the hearing advances, the interface treats documentary proof as relevant to the disputed proposition. That relevance is part of the adaptation's argument even if no character explains it. Retaining two possible verdicts cannot recover the distinction if both branches inherit the same account of the dispute.

Mitchell's discussion of ekphrasis directs attention to the addressee. A rendering can help one reader recognize the shield while enabling another to reconstruct the disposition of the figures \cite{mitchell1994ekphrasis}. A walkable hearing gives the addressee a further role. The reader can seek a better position, introduce evidence, or try to take a turn. Each operation gives an interpretation a consequence. The work then makes a claim about what can be done within the scene as well as what can be seen in it.

\section{What a cut and a corridor supply}
In Murray's rendering, the gold in the assembly is associated with the most just judgment. Place a close-up of the gold immediately after a man asserts he has paid, and it becomes easy to see that gold as the payment under dispute. The insert can supply a referent without altering the spoken line. A timeline format can preserve the exact intervals from which the image and speech came while leaving their resulting relation for a viewer to infer \cite{homerIliadMurray,otioWritingAdapter}. Opening the gold shot at its source helps identify the prop. To challenge what it stands for in the sequence, a reader also needs the placement that joined it to the line.

A spatial model can produce a comparable effect through traversal. Homer places Ocean around the shield's rim. Atlas demonstrates that separately positioned references can be made to look like places connected by a continuous passage \cite{homerIliadMurray,worldlabs2026atlas}. If an adaptation leads from the hearing to Ocean along a road, it makes a claim about the city that differs from the poem's claim about the shield. The distinction is visible in the difference between walking to a coast and pulling back from the city to see the rim. Both transitions can be good artistic choices. They give the reader different relations to remember.

The consequential unit is therefore sometimes between the visible objects. A reference to the Homeric passage on one side and to an Ocean rendering on the other does not explain who put a road between them. A reader needs to be able to cite the route as a construction, just as they need to be able to cite the placement of the gold insert.

\section{A staff and a rumor}
The elders in Murray's translation hold heralds' staves and rise to give judgment in turn \cite{homerIliadMurray}. In a world where the reader can pick up a staff, there are at least two distinct changes to stage. An object passes into a new hand, and the assembly either recognizes or rejects the holder's claim to a turn. A model that can depict both events still needs a production decision about their relation. The same initial pickup can lead to a herald yielding the floor or asking for the staff back. The difference appears in the response of others.

A rumor makes the separation more demanding. Suppose a spectator announces that the elders have awarded the gold and a shopkeeper shuts early to join the celebration. The closure can enter shared history while the reported verdict remains false. Park and colleagues demonstrate a route from natural-language memories to planning \cite{park2023generativeAgents}. That route makes the example feasible as a design problem. It does not guarantee that a later summary will keep report, belief, action, and correction apart. If the correction deletes the rumor outright, the shopkeeper's earlier action loses its grounds. If the model promotes the rumor into a fact, it invents the hearing's outcome.

These cases suggest two distinct continuity decisions. The first assigns authority to a proposition. The second records what happened because someone accepted it. A world can inherit the closure without inheriting the verdict. That separation also matters when a convincing rendering gives a disputed memory the appearance of a recorded event.

\section{A criterion for fidelity}
Fidelity here has at least two objects. Fidelity to a scene concerns the correspondence between what the source describes and what an adaptation presents. Fidelity to an inquiry concerns whether the adaptation preserves the distinctions a reader needs to challenge an interpretation. A visually convincing hearing can fail the second criterion by settling its governing question before the reader arrives. An openly fictional branch can do better if it shows what was stipulated and lets readers compare another premise.

This criterion can be tested. Produce two hearings with the same blocking and a shared receipt. Change only the translation supplied to the writers. Ask what the receipt settles. Then keep a pickup animation fixed and change only the assembly's response. Ask when readers believe they acquired authority. Finally, compare a road to Ocean with a withdrawal to the shield's rim. Ask readers to draw the relation they took the poem to establish. These tests would show where a translation, a procedural response, or a spatial transition changes the reader's account.

The analysis is limited to two English renderings of the hearing and a set of proposed adaptations. It does not adjudicate the Greek, reconstruct historical procedure, establish how a particular model responds, or report a reader study. An adaptation also has reasons to withhold its construction during a first encounter. The criterion can be met by providing a route to inspect the decision afterward.

\section{Conclusion}
The shift from imagetext to Worldtext is visible when a reader can do something the source never had to authorize. A receipt can alter the hearing, a staff can summon a response, and a path can turn the rim of an artifact into the edge of a city. Those are precise places to ask how the adaptation acquired its authority. A Worldtext reader should be able to follow such an operation back to the point where a translation, a prompt, an edit, or a continuity decision made it possible.

\begin{thebibliography}{9}
\bibitem{homerIliadMurray} Homer. \textit{Iliad}, Book 18. Translated by A. T. Murray. Theoi Classical Texts Library.
\bibitem{homerIliadButler} Homer. \textit{Iliad}, Book 18. Translated by Samuel Butler. Perseus Digital Library.
\bibitem{mitchell1994ekphrasis} Mitchell, W. J. T. 1994. Ekphrasis and the Other. In \textit{Picture Theory}. University of Chicago Press.
\bibitem{otioWritingAdapter} OpenTimelineIO. Writing an OTIO Adapter. Source Range vs Available Range.
\bibitem{worldlabs2026atlas} World Labs Team. 2026. Atlas: A World Model for Spatial Intelligence. World Labs Blog.
\bibitem{park2023generativeAgents} Park, Joon Sung, et al. 2023. Generative Agents: Interactive Simulacra of Human Behavior. arXiv:2304.03442.
\end{thebibliography}

\appendix
\section{Assembly instrument}
See \texttt{\_PROMPTS/assembly-status.txt}. It preserves an exact excerpt and records the missing full instrument.
\section{Making history}
Five Shield cards were manually transcribed. Fourteen Imagetext cards were imported from the earlier checkpoint. Codex selected and developed the argument. No human review or reader experiment has occurred. See \texttt{000\_\_MAKING\_HISTORY.txt}.
\section{Replication path}
Open \texttt{index.html} offline. The rebuild kit holds the remaining package members. The recovery script restores the index alongside them. For verification scope and commands, see\linebreak \texttt{000\_\_REBUILD.txt}.
\end{document}
